\section{What Structure Are We Approximating?}
\label{sec:target_structure}

Before introducing specific approximation methods, it is essential to clarify
\emph{what kind of structure} we aim to extract from transformer gradient matrices.
Although the gradients studied in this chapter often exhibit strong low-rank or
spiky spectra, our objective is \emph{not} to approximate an arbitrary low-rank
matrix in the classical numerical linear algebra sense.
Instead, we focus on a highly specific and structured form of low-dimensionality
that arises from shared linguistic regularities.

\paragraph{Agreement-driven structure.}
As established in Sections~\ref{sec:gradients_linear_layers} and
\ref{sec:foundational_assumptions}, the gradient matrix of a transformer linear
layer can be written as a sum of per-token outer products,
\[
G = \sum_{j,t} \delta^{(j,t)} \bigl(x^{(j,t)}\bigr)^\top.
\]
Under the assumptions that (i) grammatical structure is shared across tokens and
sequences, and (ii) semantic variation is largely independent across sequences,
the aggregated gradient naturally decomposes as
\[
G = A + R,
\]
where:
\begin{itemize}
    \item $A$ is a low-rank \emph{agreement component} induced by grammar-driven,
    globally aligned update directions; and
    \item $R$ is a high-rank residual capturing semantic and context-specific
    variation that cancels under appropriate aggregation.
\end{itemize}

The term \emph{agreement} emphasizes that $A$ arises from consistent, repeated
update signals shared across many tokens and sequences, rather than from generic
low-rank structure.

\paragraph{Why this is not generic low-rank approximation.}
In general, a low-rank matrix may have singular vectors that are sparse,
oscillatory, or orthogonal to simple probe directions.
Recovering such structure typically requires explicit singular value
decomposition or iterative orthogonalization.
In contrast, the agreement component $A$ considered here has several special
properties:
\begin{itemize}
    \item its singular vectors are dense and globally aligned,
    \item it induces non-vanishing row and column marginals,
    \item it survives simple aggregation operations such as summation or averaging.
\end{itemize}
These properties are a direct consequence of shared grammatical structure and
sequence-level semantic independence, not of low rank alone.

\paragraph{Implications for approximation.}
Because the agreement component $A$ produces coherent, non-canceling marginals,
it can be detected using extremely simple probes, such as multiplication by the
all-ones vector or by a small number of random vectors.
These probes are sufficient to recover the dominant subspace of $A$ up to scale,
without requiring explicit singular value decomposition or QR factorization.

Crucially, this approximation strategy is \emph{structure-exploiting} rather than
\emph{structure-agnostic}.
It is designed to isolate grammar-driven agreement spikes while leaving
intra-sequence semantic correlations intact, thereby preserving the information
necessary for learning semantic features.

\paragraph{Scope and limitations.}
Throughout this chapter, all approximation results should be interpreted in the
context of this agreement-driven structure.
The proposed averaging-based methods are not intended to recover arbitrary
low-rank matrices, nor to replace general-purpose SVD algorithms.
Rather, they provide a lightweight, GPU-friendly mechanism for approximating the
dominant optimization-driving structure that arises under realistic assumptions
about linguistic data and training pipelines.

With this target structure in mind, we now turn to averaging operations as simple
yet effective probes for isolating agreement-driven components in gradient
matrices.